\documentclass{ximera}

\input{../../preamble.tex}

\title{Shifted Exponential}

\begin{document}

\begin{abstract}
almost exponential
\end{abstract}
\maketitle





\section*{Shifted Exponential Functions}


Shifted Exponential functions are exponential functions with a nonzero constant term.

They exhibit all of the characteristics of exponential functions, if you take away the constant term.


\begin{definition} \textbf{\textcolor{green!50!black}{Shifted Exponential Functions}}

Shifted exponential functions are exponential functions with a number added on.  They \textbf{\textcolor{purple!85!blue}{can}} be written in the form


\[      f(x) = A \cdot r^{B \, x + C} + D   \]

where $A$, $B$, $C$, and $D$ are real numbers, $A \ne 0$ and $B \ne 0$ and $D \ne 0$, and $r$ is a positive real number.


\end{definition}

\textbf{Note:}  In the template for shifted exponential functions, There is a leading coefficient for the function and there is a leading coefficient for the linear function inside the exponent. 


The main different between shifted exponential and exponential functions is that while exponential functions do not have zeros, a shifted exponential function may have a zero. 







Shifted exponential functions can be thought of as compositions our our Core exponential funcitons with linear functions.



\begin{idea} \textbf{\textcolor{blue!55!black}{Shifted Exponential Functions}}


Let $E(k) = r^k $ be a Core exponential function.

Let $L_{in}(t) = B \, t + C$ be a linear function.

Let $L_{out}(y) = A \, y + D$ be a linear function.


General shifted exponential functions are constructed by composing these component functions.


\[
(L_{out} \circ E \circ L_{in})(x) = A \, r^{B \,x + C} + D
\]




\end{idea}




\begin{example}

Here is the graph of $g(x) = 2^x - 3$.

\begin{image}
\begin{tikzpicture} 
  \begin{axis}[
            domain=-10:10, ymax=10, xmax=10, ymin=-10, xmin=-10,
            axis lines =center, xlabel=$x$, ylabel=$y$, 
            ytick={-10,-8,-6,-4,-2,2,4,6,8,10},
            xtick={-10,-8,-6,-4,-2,2,4,6,8,10},
            ticklabel style={font=\scriptsize},
            every axis y label/.style={at=(current axis.above origin),anchor=south},
            every axis x label/.style={at=(current axis.right of origin),anchor=west},
            axis on top
          ]
          
          \addplot [line width=2, penColor, smooth,samples=200,domain=(-9:3.2),<->] {2^x-3};
          \addplot [line width=1, gray, dashed,domain=(-9:9),<->] ({x},{-3});

           

  \end{axis}
\end{tikzpicture}
\end{image}


$\log_2(3)$ is the zero of $g(x)$.



\end{example}










\begin{center}
\textbf{\textcolor{green!50!black}{ooooo-=-=-=-ooOoo-=-=-=-ooooo}} \\

more examples can be found by following this link\\ \link[More Examples of Elementary Functions]{https://ximera.osu.edu/csccmathematics/precalculus1/precalculus1/elementaryLibrary2/examples/exampleList}

\end{center}




\end{document}
